\documentclass[11pt]{article}

\usepackage[margin=1in]{geometry}
\usepackage{amsmath, amssymb, amsthm, mathtools}
\usepackage{booktabs}
\usepackage{enumitem}
\usepackage{xcolor}
\usepackage[colorlinks=true, linkcolor=blue, citecolor=blue, urlcolor=blue]{hyperref}

\title{Self-Modulated Neural Networks via Sigma-Algebra Edge Addressing}
\author{}
\date{}

\newtheorem{definition}{Definition}
\newtheorem{proposition}{Proposition}
\newtheorem{remark}{Remark}

\newcommand{\R}{\mathbb{R}}
\newcommand{\1}{\mathbf{1}}

\begin{document}

\maketitle

\begin{abstract}
We propose a simple framework for self-modulated neural networks. A neural network is equipped with two types of output heads: ordinary task heads, denoted by \(y\), and internal modulator heads, denoted by \(q\). The task heads are trained by the usual task loss and ordinary backpropagation, while the modulator heads generate internal signals that regulate the network's own edge weights. The central mathematical observation is that \(m\) binary generators can distinguish up to \(2^m\) controllable edge weights through the sigma-algebra they generate. Thus, a logarithmic number of modulator heads can provide distinct addresses for many edge weights. This addressability should not be confused with full independent controllability: \(O(\log N)\) heads cannot assign arbitrary independent updates to \(N\) weights in one step. Instead, they provide a compact, energy-normalized modulation mechanism whose reachable linear signal space has dimension at most \(m\). We formulate the model, clarify what is and is not claimed, and present a hand-computable example on a small multilayer perceptron.
\end{abstract}

\section{Sigma-Algebra Addressability of Edge Weights}

\subsection{Edge weights as a finite measurable space}

Let
\[
E=\{e_1,\dots,e_N\}
\]
denote the set of controllable edge weights in a neural network. Here \(N\) is not necessarily the total number of graph edges in the computational graph. Rather, \(N\) is the number of trainable edge weights that we choose to expose to the internal modulation mechanism.

For example, in a residual block of the form
\[
x \mapsto x+F_\theta(x),
\]
the identity skip connection \(x\mapsto x\) is typically a fixed structural edge. It does not have a trainable scalar weight and therefore should not be counted in \(E\). If, however, the skip connection has a trainable gate,
\[
x \mapsto \alpha x+F_\theta(x),
\]
then the scalar \(\alpha\) may be included in \(E\).

We treat \(E\) as a finite measurable space. The full discrete sigma-algebra on \(E\) is the power set
\[
\mathcal P(E).
\]
In this setting, distinguishing all controllable edge weights means generating a sigma-algebra whose atoms are singletons \(\{e_j\}\).

\subsection{Minimal number of heads for distinguishing \texorpdfstring{$N$}{N} edges}

Suppose we have \(m\) modulator heads
\[
q_1,\dots,q_m.
\]
For the purpose of addressability, each modulator head is associated with a subset
\[
S_i\subseteq E.
\]
The membership of an edge \(e_j\) in these sets defines its binary address
\[
a_j
=
\big(
\1_{e_j\in S_1},
\dots,
\1_{e_j\in S_m}
\big)
\in \{0,1\}^m.
\]
If two edges have the same address, then no set generated by \(S_1,\dots,S_m\) can separate them. If all addresses are distinct, then every edge can be identified as an atom of the generated sigma-algebra.

\begin{proposition}[Minimal addressability]
Let \(E=\{e_1,\dots,e_N\}\). If the all-zero address is allowed, then the minimum number of binary generators needed to distinguish all \(N\) elements is
\[
m_{\min}=\lceil \log_2 N\rceil.
\]
If, in addition, every edge must be directly reached by at least one head, so that the all-zero address is forbidden, then the minimum number is
\[
m_{\min}^{\mathrm{covered}}
=
\lceil \log_2(N+1)\rceil.
\]
\end{proposition}

\begin{proof}
With \(m\) binary generators, there are at most \(2^m\) distinct binary addresses. Hence distinguishing \(N\) elements requires
\[
N\leq 2^m,
\]
which gives
\[
m\geq \lceil \log_2 N\rceil.
\]
This bound is attainable by assigning distinct binary codes to the \(N\) elements.

If the all-zero address is forbidden, then only \(2^m-1\) nonzero addresses are available. Thus we require
\[
N\leq 2^m-1,
\]
equivalently
\[
N+1\leq 2^m.
\]
Therefore
\[
m\geq \lceil \log_2(N+1)\rceil.
\]
This bound is again attainable by assigning distinct nonzero binary codes.
\end{proof}

The generated sigma-algebra is
\[
\sigma(S_1,\dots,S_m).
\]
If the address of \(e_j\) is \(a_j=(a_{j1},\dots,a_{jm})\), then its atom is
\[
\{e_j\}
=
\bigcap_{i:a_{ji}=1}S_i
\cap
\bigcap_{i:a_{ji}=0}S_i^c.
\]
Thus the sigma-algebra generated by \(m\) heads can distinguish all \(N\) edges whenever the addresses are distinct.

\begin{remark}
In the rest of the report, we use \(m\) to denote a logarithmic number of modulator heads. If direct coverage is required, one should choose
\[
m=\lceil \log_2(N+1)\rceil.
\]
For large \(N\), this is essentially the same logarithmic scaling as \(\lceil \log_2 N\rceil\).
\end{remark}

\begin{proposition}[Linear control dimension bound]
Let
\[
s=\widetilde B q,
\qquad
\widetilde B\in\R^{N\times m},
\qquad
q\in\R^m.
\]
Then the set of linear modulation signals reachable in one step is
\[
\mathcal S
=
\{\widetilde B q:q\in\R^m\}
=
\operatorname{im}(\widetilde B),
\]
which is a linear subspace of \(\R^N\) of dimension
\[
\dim(\mathcal S)=\operatorname{rank}(\widetilde B)\le m.
\]
In particular, if \(N>m\), then one cannot in general realize an arbitrary edge-wise signal in \(\R^N\) by choosing \(q\).
\end{proposition}

\begin{proof}
The reachable set is precisely the image of the linear map \(q\mapsto \widetilde B q\). Its dimension is therefore the rank of \(\widetilde B\), which is at most \(\min(N,m)\), hence at most \(m\). If \(N>m\), then \(\operatorname{rank}(\widetilde B)\le m<N\), so \(\operatorname{im}(\widetilde B)\) is a proper subspace of \(\R^N\). Therefore not every vector in \(\R^N\) is reachable.
\end{proof}

\subsection{Distinguishability versus independent controllability}

The previous argument concerns distinguishability, not full independent controllability. This distinction is essential.

Let \(w_j\) denote the weight associated with edge \(e_j\). Full independent controllability would mean that in one step we can prescribe an arbitrary update vector
\[
(\Delta w_1,\dots,\Delta w_N)\in \R^N.
\]
In general, achieving this requires \(N\) independent control directions.

By contrast, \(m=O(\log N)\) modulator heads provide a compact addressing system. They allow each edge to have a distinct response pattern to the modulator signals, but they do not allow arbitrary independent updates of all weights in one step.

Thus our goal is not
\[
\text{full independent control of every edge weight}.
\]
Rather, our goal is
\[
\text{addressable internal modulation of all controllable edge weights}.
\]
In short,
\[
O(\log N)
\]
is an addressing dimension, not a full control dimension.

\section{Self-Modulated Neural Network Architecture}

\subsection{Task heads and modulator heads}

We consider a neural network \(A_\theta\) whose output is split into two parts:
\[
A_\theta(x)
=
\big(
y_\theta(x),q_\theta(x)
\big).
\]
The vector
\[
y_\theta(x)
\]
contains ordinary task heads. These are used for the external task, such as classification, regression, or prediction. The vector
\[
q_\theta(x)
=
(q_1(x),\dots,q_m(x))
\]
contains modulator heads. These heads are used to produce internal signals that regulate the network's own edge weights.

We assume
\[
q_\theta(x)\in[-1,1]^m.
\]
For example, one may compute raw modulator outputs \(q_{\mathrm{raw}}(x)\in\R^m\) and set
\[
q_\theta(x)=\tanh(q_{\mathrm{raw}}(x)).
\]

Given a training sample \((x_t,y_t^\ast)\), the task loss is
\[
L_t=\ell(y_\theta(x_t),y_t^\ast).
\]
The ordinary backpropagation update for an edge weight \(w_j\) is
\[
\Delta w_j^{\mathrm{bp}}(t)
=
-\eta_{\mathrm{bp}}
\frac{\partial L_t}{\partial w_j}.
\]
This is the standard task-driven update.

The modulator heads generate a separate internal update. The internal update is not defined by the gradient. Instead, it is produced by the modulator outputs \(q_\theta(x_t)\), the head-to-edge assignment, and an energy normalization rule.

\begin{remark}[Current experimental regime]
For the first version of this project, we fix the assignment matrix \(B\) once at initialization and keep it static throughout training. Thus the set of synapses influenced by each modulator head is fixed in advance.

We also do not introduce a separate auxiliary loss or outer-loop objective specifically for the \(q\)-heads. Instead, training starts in a regime where ordinary backpropagation dominates, i.e. \(\lambda_t\) is small, and then gradually transitions toward stronger internal modulation by increasing \(\lambda_t\).

The working hypothesis is not that the \(q\)-heads are explicitly taught a modulation rule. Rather, the experiment asks whether a useful self-modulation pattern can emerge gradually under this schedule once the network has first been shaped by ordinary task-driven learning.
\end{remark}

\subsection{Randomized but constrained head-to-edge assignment}

Let
\[
E=\{e_1,\dots,e_N\}
\]
be the controllable edge set. We assign the \(m\) modulator heads to the \(N\) edges using an incidence matrix
\[
B\in\{0,1\}^{N\times m}.
\]
The entry \(B_{ji}\) is defined by
\[
B_{ji}=1
\quad\Longleftrightarrow\quad
e_j\in S_i,
\]
where \(S_i\subseteq E\) is the set of edges controlled by head \(q_i\).

The \(j\)-th row
\[
B_{j,:}
\]
is the address of edge \(e_j\). The assignment matrix may be sampled or designed at initialization, but in the current experiment it is thereafter kept fixed. The constraints are as follows.

Under the full-coverage requirement \(B_{j,:}\neq 0\), feasibility already requires
\[
N\le 2^m-1,
\]
which is equivalent to
\[
m\ge \lceil \log_2(N+1)\rceil.
\]

First, we require full coverage:
\[
B_{j,:}\neq 0,
\qquad
j=1,\dots,N.
\]
This means that every controllable edge receives at least one modulator signal.

Second, we require distinguishability:
\[
B_{j,:}\neq B_{\ell,:},
\qquad
j\neq \ell.
\]
This means that no two controllable edges have the same address.

Third, we require balance and dispersion. Define
\[
k_i=\sum_{j=1}^N B_{ji}.
\]
Thus \(k_i\) is the number of edges controlled by head \(q_i\). We do not want some heads to control almost all edges while others control almost none. Hence the distribution of \(k_i\) should be reasonably balanced.

Moreover, if the edges are arranged by layers,
\[
E=E^{(1)}\cup\cdots\cup E^{(L)},
\]
then each \(S_i\) should be dispersed across the network rather than concentrated in a single layer or local block. A natural dispersion target is
\[
\frac{|S_i\cap E^{(\ell)}|}{|E^{(\ell)}|}
\approx
\frac{k_i}{N}.
\]
This says that the fraction of layer-\(\ell\) edges controlled by head \(q_i\) should roughly match the global fraction of edges controlled by that head.

Thus the assignment is random only after enforcing coverage, distinguishability, balance, and dispersion.

\subsection{Energy-normalized internal modulation}

Each modulator head has finite energy. If head \(q_i\) controls \(k_i\) edges, then its signal should be diluted across these \(k_i\) edges. We adopt an \(L^2\) energy normalization:
\[
\text{amplitude per controlled edge}
=
\frac{1}{\sqrt{k_i}}.
\]
This normalization ensures that if head \(q_i\) distributes the same magnitude to all controlled edges, then the squared energy of that head remains \(O(1)\):
\[
\sum_{j:B_{ji}=1}
\left(
\frac{q_i}{\sqrt{k_i}}
\right)^2
=
q_i^2.
\]

The internal modulation signal received by edge \(e_j\) is
\[
s_j(t)
=
\sum_{i=1}^m
\frac{B_{ji}}{\sqrt{k_i}}q_i(t).
\]
Here
\[
q_i(t):=q_i(x_t),
\]
so the modulation signal is produced from the current network response to the current input.
Equivalently, define the normalized incidence matrix
\[
\widetilde B_{ji}
=
\frac{B_{ji}}{\sqrt{k_i}}.
\]
Then
\[
s(t)=\widetilde B q(t).
\]

We define the internal modulation update by
\[
\Delta w_j^{\mathrm{int}}(t)
=
\eta_{\mathrm{int}}
\tanh(\gamma s_j(t)),
\]
where \(\eta_{\mathrm{int}}>0\) is the internal modulation step size and \(\gamma>0\) controls the sharpness of the modulation response.

The total update is the mixture
\[
w_j(t+1)
=
w_j(t)
+
(1-\lambda_t)\Delta w_j^{\mathrm{bp}}(t)
+
\lambda_t\Delta w_j^{\mathrm{int}}(t),
\]
where
\[
\lambda_t\in[0,1].
\]
At the beginning of training, one may choose
\[
\lambda_t\approx 0,
\]
so that ordinary backpropagation dominates. Later, one may gradually increase \(\lambda_t\), allowing the internal modulation mechanism to play a larger role.

The crucial point is that the internal modulation update does not directly use the gradient:
\[
\Delta w_j^{\mathrm{int}}(t)
=
\eta_{\mathrm{int}}\tanh(\gamma s_j(t)).
\]
The gradient appears only in the ordinary backpropagation term
\[
\Delta w_j^{\mathrm{bp}}(t)
=
-\eta_{\mathrm{bp}}
\frac{\partial L_t}{\partial w_j}.
\]

\section{A Hand-Computable MLP Example and Algorithmic Template}

\subsection{A small MLP with \texorpdfstring{$y$}{y}-heads and \texorpdfstring{$q$}{q}-heads}

We now give a hand-computable example. Let
\[
x=(x_1,x_2)\in\R^2.
\]
Consider a one-hidden-layer MLP with three hidden neurons
\[
h_1,h_2,h_3.
\]
To keep the calculation transparent, assume that the hidden activation is the identity function, or equivalently that all ReLU units are in their active region.

Let the input-to-hidden weights be
\[
w_1,\dots,w_6.
\]
Define
\[
h_1=w_1x_1+w_2x_2,
\]
\[
h_2=w_3x_1+w_4x_2,
\]
\[
h_3=w_5x_1+w_6x_2.
\]
Thus there are
\[
N=6
\]
controllable input-to-hidden edges. Since
\[
\lceil \log_2(6+1)\rceil=3,
\]
we use
\[
m=3
\]
modulator heads.

The task output is
\[
y=v_1h_1+v_2h_2+v_3h_3.
\]
The modulator output is
\[
q=(q_1,q_2,q_3)\in[-1,1]^3.
\]
Therefore the network output is
\[
A_\theta(x)=(y,q).
\]
In the hand calculation below, \(q\) is treated as a given numerical output. In an actual implementation, \(q\) would be produced by the network's modulator heads.

\subsection{Sigma-algebra assignment and edge-level modulation signals}

We assign the six controllable edges to the three modulator heads using the incidence matrix
\[
B=
\begin{pmatrix}
0&0&1\\
0&1&0\\
0&1&1\\
1&0&0\\
1&0&1\\
1&1&0
\end{pmatrix}.
\]
Equivalently,
\[
\begin{array}{c|ccc}
\toprule
\text{edge} & q_1 & q_2 & q_3\\
\midrule
e_1=w_1 & 0 & 0 & 1\\
e_2=w_2 & 0 & 1 & 0\\
e_3=w_3 & 0 & 1 & 1\\
e_4=w_4 & 1 & 0 & 0\\
e_5=w_5 & 1 & 0 & 1\\
e_6=w_6 & 1 & 1 & 0\\
\bottomrule
\end{array}
\]
The corresponding controlled sets are
\[
S_1=\{e_4,e_5,e_6\},
\]
\[
S_2=\{e_2,e_3,e_6\},
\]
\[
S_3=\{e_1,e_3,e_5\}.
\]
Each head controls three edges:
\[
k_1=k_2=k_3=3.
\]
Thus the energy normalization factor is
\[
\frac{1}{\sqrt{3}}.
\]

The edge-level modulation signals are
\[
s_j=\sum_{i=1}^3\frac{B_{ji}}{\sqrt{3}}q_i.
\]
Expanding this gives
\[
s_1=\frac{q_3}{\sqrt3},
\]
\[
s_2=\frac{q_2}{\sqrt3},
\]
\[
s_3=\frac{q_2+q_3}{\sqrt3},
\]
\[
s_4=\frac{q_1}{\sqrt3},
\]
\[
s_5=\frac{q_1+q_3}{\sqrt3},
\]
\[
s_6=\frac{q_1+q_2}{\sqrt3}.
\]
Thus three modulator heads generate six distinct edge-level modulation signals. However, these signals are still functions of the same three-dimensional vector \(q\). Therefore this construction provides addressable modulation, not arbitrary independent six-dimensional control.

\subsection{One-step mixed update: backprop plus internal modulation}

We now compute one training step explicitly. Take
\[
x=(1,2),
\qquad
y^\ast=1.
\]
Let
\[
w_1=0.1,\quad w_2=0.2,
\]
\[
w_3=-0.1,\quad w_4=0.1,
\]
\[
w_5=0.05,\quad w_6=-0.05.
\]
Let the output weights be
\[
v=(v_1,v_2,v_3)=(1,-1,0.5).
\]
Then
\[
h_1=0.1\cdot 1+0.2\cdot 2=0.5,
\]
\[
h_2=-0.1\cdot 1+0.1\cdot 2=0.1,
\]
\[
h_3=0.05\cdot 1-0.05\cdot 2=-0.05.
\]
Therefore
\[
y=h_1-h_2+0.5h_3
=
0.5-0.1-0.025
=
0.375.
\]
Use the squared loss
\[
L=\frac12(y-y^\ast)^2.
\]
The prediction error is
\[
\delta=y-y^\ast=0.375-1=-0.625.
\]

For the input-to-hidden weights, the ordinary backpropagation gradient is
\[
\frac{\partial L}{\partial w_{r,a}}
=
\delta v_r x_a,
\]
where \(w_{r,a}\) denotes the edge from input coordinate \(x_a\) to hidden unit \(h_r\). Therefore
\[
g^{\mathrm{bp}}
=
\left(
-0.625,\,
-1.25,\,
0.625,\,
1.25,\,
-0.3125,\,
-0.625
\right).
\]
If
\[
\eta_{\mathrm{bp}}=0.01,
\]
then
\[
\Delta w^{\mathrm{bp}}
=
-\eta_{\mathrm{bp}}g^{\mathrm{bp}}
=
\left(
0.00625,\,
0.0125,\,
-0.00625,\,
-0.0125,\,
0.003125,\,
0.00625
\right).
\]

Now suppose the modulator heads output
\[
q=(0.6,-0.2,0.8).
\]
Using the expressions from the previous subsection,
\[
s_1=\frac{0.8}{\sqrt3}\approx0.462,
\]
\[
s_2=\frac{-0.2}{\sqrt3}\approx-0.115,
\]
\[
s_3=\frac{-0.2+0.8}{\sqrt3}\approx0.346,
\]
\[
s_4=\frac{0.6}{\sqrt3}\approx0.346,
\]
\[
s_5=\frac{0.6+0.8}{\sqrt3}\approx0.808,
\]
\[
s_6=\frac{0.6-0.2}{\sqrt3}\approx0.231.
\]

Let
\[
\eta_{\mathrm{int}}=0.01,
\qquad
\gamma=1.
\]
The internal modulation update is
\[
\Delta w_j^{\mathrm{int}}
=
\eta_{\mathrm{int}}\tanh(s_j).
\]
Since
\[
\tanh(s)
\approx
(0.432,\,-0.115,\,0.333,\,0.333,\,0.669,\,0.227),
\]
we obtain
\[
\Delta w^{\mathrm{int}}
\approx
\left(
0.00432,\,
-0.00115,\,
0.00333,\,
0.00333,\,
0.00669,\,
0.00227
\right).
\]

Finally, choose
\[
\lambda=0.2.
\]
The mixed update is
\[
\Delta w
=
(1-\lambda)\Delta w^{\mathrm{bp}}
+
\lambda \Delta w^{\mathrm{int}}.
\]
Thus
\[
\Delta w
\approx
0.8\Delta w^{\mathrm{bp}}
+
0.2\Delta w^{\mathrm{int}},
\]
which gives
\[
\Delta w
\approx
\left(
0.00586,\,
0.00977,\,
-0.00433,\,
-0.00933,\,
0.00384,\,
0.00545
\right).
\]
Therefore the updated input-to-hidden weights are
\[
w_j^+
=
w_j+\Delta w_j,
\qquad
j=1,\dots,6.
\]

This example shows the complete mechanism. The \(y\)-head produces the task output and the ordinary backpropagation update. The \(q\)-heads produce internal modulation signals. The incidence matrix \(B\) maps three modulator heads to six controllable edges. The factor \(1/\sqrt{k_i}\) implements energy-normalized signal distribution. The parameter \(\lambda\) controls the mixture between ordinary backpropagation and internal self-modulation.

\subsection{Minimal algorithmic template}

The full training mechanism can be summarized at the level of one sample \((x_t,y_t^\ast)\) as follows:
\begin{enumerate}[leftmargin=*]
\item Choose a controllable edge set \(E\), a number of modulator heads \(m\), and an assignment matrix \(B\) satisfying coverage and distinguishability.
\item Run a forward pass to obtain
\[
(y_t,q_t)=A_\theta(x_t).
\]
\item Compute the task loss
\[
L_t=\ell(y_t,y_t^\ast)
\]
and the ordinary task-driven update
\[
\Delta w^{\mathrm{bp}}(t)=-\eta_{\mathrm{bp}}\nabla_w L_t.
\]
\item Form the edge-level modulation signal
\[
s(t)=\widetilde B q_t
\]
and the internal update
\[
\Delta w^{\mathrm{int}}_j(t)=\eta_{\mathrm{int}}\tanh(\gamma s_j(t)).
\]
\item Mix the two updates:
\[
w(t+1)=w(t)+(1-\lambda_t)\Delta w^{\mathrm{bp}}(t)+\lambda_t\Delta w^{\mathrm{int}}(t).
\]
\item Update the remaining non-modulated parameters according to the chosen learning regime. If the \(q\)-heads are trainable, their own parameters must be trained by one of the regimes discussed earlier.
\end{enumerate}

This note establishes the addressing and update mechanism. It does not yet claim that a particular learning rule for the \(q\)-heads is optimal.

\section*{Summary of the Core Mechanism}

The proposed framework can be summarized by the following equations:
\[
A_\theta(x)=(y_\theta(x),q_\theta(x)),
\]
\[
q_\theta(x)\in[-1,1]^m,
\]
\[
m=O(\log N),
\]
\[
s_j(t)
=
\sum_{i=1}^m
\frac{B_{ji}}{\sqrt{k_i}}q_i(t),
\]
\[
\Delta w_j^{\mathrm{int}}(t)
=
\eta_{\mathrm{int}}\tanh(\gamma s_j(t)),
\]
\[
\Delta w_j^{\mathrm{bp}}(t)
=
-\eta_{\mathrm{bp}}
\frac{\partial L_t}{\partial w_j},
\]
and
\[
w_j(t+1)
=
w_j(t)
+
(1-\lambda_t)\Delta w_j^{\mathrm{bp}}(t)
+
\lambda_t\Delta w_j^{\mathrm{int}}(t).
\]

The central claim is that logarithmically many modulator heads can provide addressability for many edge weights through a finite sigma-algebra structure. This does not imply full independent controllability. Instead, it yields an energy-normalized internal modulation mechanism that can be combined with ordinary backpropagation during training.

\end{document}
